% GENERATED FILE. Edit canonical lessons, then run npm run generate:books.
% canonical-source-hash: fnv1a64:fc0ea36cede55654
% canonical-lessons: FR-C309-maternite, FR-C309-bapteme, FR-C309-baptiser, FR-C309-poussette, FR-C309-enfance

\chapter{Les premières années}
\label{ch:fr-les-premieres-annees}
\hlchaptermodality{309}

\begin{chapteropening}
\textbf{By the end of this chapter:} \emph{I can talk about early childhood: the maternity ward, a christening, christening a child, a pushchair and childhood.}

Complete the last lesson of chapter 309: I can talk about early childhood: the maternity ward, a christening, christening a child, a pushchair and childhood.
\end{chapteropening}

\section[la maternité]{la maternité — the maternity hospital; motherhood}

\label{lesson:FR-C309-maternite}



\begin{quote}
Before the new one: say the French for to breastfeed, then the French for the cradle.
\end{quote}

\subsection*{You'll want to know: la maternité}
\textbf{la maternité} — \textquotedblleft{}the maternity hospital; motherhood\textquotedblright{}.

\begin{quote}
Le bébé est né à la maternité de la ville.
\end{quote}

\begin{quote}
\emph{The baby was born at the town's maternity hospital.}
\end{quote}

\emph{Le congé de maternité} — maternity leave.

\begin{cousinweb}
\emph{maternité} is a learned borrowing from Latin \emph{māternitās} \textquotedblleft{}motherhood\textquotedblright{} (EtymDB), built on \emph{māter} \textquotedblleft{}mother\textquotedblright{}.
\end{cousinweb}

\begin{grammarlens}[title={What you've built}]
One more word for telling, explaining and describing at B1.
\end{grammarlens}

\subsection*{Your turn}
\begin{itemize}
\raggedright
  \item \emph{Say it:} \emph{la maternité}
  \item \emph{Say it:} \emph{la maternité}, once more
  \item \emph{Say it:} a sentence of your own with \emph{la maternité}
  \item \emph{Recall:} say the French for to breastfeed, then the French for the cradle, then say \emph{la maternité} again
\end{itemize}

\begin{tcolorbox}[breakable,colback=teal!4,colframe=teal!35!black,title={Before you move on}]
What does la maternité mean? (\textquotedblleft{}The maternity hospital; motherhood\textquotedblright{}.) Say it once more.
\end{tcolorbox}
\section[le baptême]{le baptême — the christening, the baptism}

\label{lesson:FR-C309-bapteme}



\begin{quote}
Before the new one: say the French for the cradle, then the French for the maternity hospital.
\end{quote}

\subsection*{You'll want to know: le baptême}
\textbf{le baptême} — \textquotedblleft{}the christening, the baptism\textquotedblright{}.

\begin{quote}
Toute la famille est venue au baptême.
\end{quote}

\begin{quote}
\emph{The whole family came to the christening.}
\end{quote}

The \emph{p} is silent. \emph{Le baptême de l'air} — a first flight.

\begin{cousinweb}
\emph{baptême} comes through Old French \emph{batesme} (EtymDB) from Church Latin \emph{baptisma}, a borrowing from Greek \emph{báptisma} \textquotedblleft{}a dipping\textquotedblright{}.
\end{cousinweb}

\begin{grammarlens}[title={What you've built}]
One more word for telling, explaining and describing at B1.
\end{grammarlens}

\subsection*{Your turn}
\begin{itemize}
\raggedright
  \item \emph{Say it:} \emph{le baptême}
  \item \emph{Say it:} \emph{le baptême}, once more
  \item \emph{Say it:} a sentence of your own with \emph{le baptême}
  \item \emph{Recall:} say the French for the cradle, then the French for the maternity hospital, then say \emph{le baptême} again
\end{itemize}

\begin{tcolorbox}[breakable,colback=teal!4,colframe=teal!35!black,title={Before you move on}]
What does le baptême mean? (\textquotedblleft{}The christening, the baptism\textquotedblright{}.) Say it once more.
\end{tcolorbox}
\section[baptiser]{baptiser — to christen, to baptise; to name}

\label{lesson:FR-C309-baptiser}



\begin{quote}
Before the new one: say the French for the maternity hospital, then the French for the christening.
\end{quote}

\subsection*{You'll want to know: baptiser}
\textbf{baptiser} — \textquotedblleft{}to christen, to baptise; to name\textquotedblright{}.

\begin{quote}
Le prêtre a baptisé l'enfant dimanche.
\end{quote}

\begin{quote}
\emph{The priest christened the child on Sunday.}
\end{quote}

Also to give a name: \emph{ils ont baptisé leur bateau}.

\begin{cousinweb}
\emph{baptiser} comes through Old French \emph{baptisier} (EtymDB) from Church Latin \emph{baptizāre}, a borrowing from Greek \emph{baptízein} \textquotedblleft{}to dip\textquotedblright{}.
\end{cousinweb}

\begin{grammarlens}[title={What you've built}]
One more word for telling, explaining and describing at B1.
\end{grammarlens}

\subsection*{Your turn}
\begin{itemize}
\raggedright
  \item \emph{Say it:} \emph{baptiser}
  \item \emph{Say it:} \emph{baptiser}, once more
  \item \emph{Say it:} a sentence of your own with \emph{baptiser}
  \item \emph{Recall:} say the French for the maternity hospital, then the French for the christening, then say \emph{baptiser} again
\end{itemize}

\begin{tcolorbox}[breakable,colback=teal!4,colframe=teal!35!black,title={Before you move on}]
What does baptiser mean? (\textquotedblleft{}To christen, to baptise; to name\textquotedblright{}.) Say it once more.
\end{tcolorbox}
\section[la poussette]{la poussette — the pushchair, the stroller}

\label{lesson:FR-C309-poussette}



\begin{quote}
Before the new one: say the French for the christening, then the French for to christen.
\end{quote}

\subsection*{You'll want to know: la poussette}
\textbf{la poussette} — \textquotedblleft{}the pushchair, the stroller\textquotedblright{}.

\begin{quote}
Elle se promène au parc avec la poussette.
\end{quote}

\begin{quote}
\emph{She walks in the park with the pushchair.}
\end{quote}

Built on \emph{pousser}, to push.

\begin{cousinweb}
\emph{poussette} is built on \emph{pousser} \textquotedblleft{}to push\textquotedblright{}, which is inherited from Latin \emph{pulsāre} \textquotedblleft{}to beat, to push\textquotedblright{} (Lewis \& Short).
\end{cousinweb}

\begin{grammarlens}[title={What you've built}]
One more word for telling, explaining and describing at B1.
\end{grammarlens}

\subsection*{Your turn}
\begin{itemize}
\raggedright
  \item \emph{Say it:} \emph{la poussette}
  \item \emph{Say it:} \emph{la poussette}, once more
  \item \emph{Say it:} a sentence of your own with \emph{la poussette}
  \item \emph{Recall:} say the French for the christening, then the French for to christen, then say \emph{la poussette} again
\end{itemize}

\begin{tcolorbox}[breakable,colback=teal!4,colframe=teal!35!black,title={Before you move on}]
What does la poussette mean? (\textquotedblleft{}The pushchair, the stroller\textquotedblright{}.) Say it once more.
\end{tcolorbox}
\section[l'enfance]{l'enfance — childhood}

\label{lesson:FR-C309-enfance}



\begin{quote}
Before the new one: say the French for to christen, then the French for the pushchair.
\end{quote}

\subsection*{You'll want to know: l'enfance}
\textbf{l'enfance} — \textquotedblleft{}childhood\textquotedblright{}.

\begin{quote}
Il a passé son enfance à la campagne.
\end{quote}

\begin{quote}
\emph{He spent his childhood in the countryside.}
\end{quote}

\emph{Un ami d'enfance} — a childhood friend.

\begin{cousinweb}
\emph{enfance} comes from Old French \emph{enfance}, from Latin \emph{īnfantia} \textquotedblleft{}childhood\textquotedblright{} (EtymDB), built on \emph{īnfāns} \textquotedblleft{}unable to speak\textquotedblright{}.
\end{cousinweb}

\begin{grammarlens}[title={What you've built}]
One more word for telling, explaining and describing at B1.
\end{grammarlens}

\subsection*{Your turn}
\begin{itemize}
\raggedright
  \item \emph{Say it:} \emph{l'enfance}
  \item \emph{Say it:} \emph{l'enfance}, once more
  \item \emph{Say it:} a sentence of your own with \emph{l'enfance}
  \item \emph{Recall:} say the French for to christen, then the French for the pushchair, then say \emph{l'enfance} again
\end{itemize}

\begin{tcolorbox}[breakable,colback=teal!4,colframe=teal!35!black,title={Before you move on}]
What does l'enfance mean? (\textquotedblleft{}Childhood\textquotedblright{}.) Say it once more.
\end{tcolorbox}
